\documentclass{article}

% packages for math
\usepackage{amsthm, amsmath, amssymb, amsfonts}

% package for including images
\usepackage{graphicx}

% environment for solutions
\theoremstyle{remark} \newtheorem*{solution}{Solution}

% capital letters for problem parts
\renewcommand{\theenumi}{\Alph{enumi}}

% no page numbers
\pagenumbering{gobble}

% codeblock
\usepackage{listings}
\usepackage{xcolor}

\definecolor{codegreen}{rgb}{0,0.6,0}
\definecolor{codegray}{rgb}{0.5,0.5,0.5}
\definecolor{codepurple}{rgb}{0.58,0,0.82}
\definecolor{backcolour}{rgb}{0.95,0.95,0.95}

\lstdefinestyle{mystyle}{
    backgroundcolor=\color{backcolour},
    commentstyle=\color{codegreen},
    keywordstyle=[1]{\color{blue}},
    keywordstyle=[2]{\color{blue}},
    numberstyle=\tiny\color{codegray},
    stringstyle=\color{codepurple},
    basicstyle=\ttfamily,
    breakatwhitespace=false,
    breaklines=true,
    captionpos=b,
    keepspaces=true,
    showspaces=false,
    showstringspaces=false,
    showtabs=false,
    tabsize=2
}

\lstset{style=mystyle}

\lstdefinelanguage{ocaml}{
  morekeywords=[1]{let,in,match,with,rec,and},
  sensitive=true,
  morecomment=[s]{(*}{*)},
  morestring=[b]",
}


\title{Practice Midterm}
\author{CAS CS 320: Principles of Programming Languages}

\begin{document}
\maketitle

\noindent Name:

\bigskip

\noindent BUID:

\bigskip

\noindent Location:

\bigskip

\begin{itemize}
\item You will have approximately 75 minutes to complete this exam.
\item Make sure to read every question, some are easier than others.
\item Please write your name and BUID \textbf{on every page}.
\end{itemize}

\pagebreak
\textit{(Extra page)}

\pagebreak
\section{Shadowing}

Consider the following definition.
\begin{lstlisting}[language=ocaml]
  let x y =
    let x = y in
    let y x = x + 1 in
    let x y = y x in
    x y
\end{lstlisting}
\begin{enumerate}
\item
  What is the the type of \texttt{x}?
  \vspace{1in}
\item
  What is the value of \texttt{x 2}?
\end{enumerate}

\pagebreak
\textit{(Extra page)}

\pagebreak
\section{Number of Digits}

Implement the function \texttt{num\_digits} which, given an integer $n$, returns the number of digits in $n$.
You may only use arithmetic operations.
\bigskip

\noindent
\texttt{let num\_digits (n : int) : int =}
\vfill

\noindent
\texttt{let \_ = assert (num\_digits 12345 = 5)}

\noindent
\texttt{let \_ = assert (num\_digits -10 = 2)}

\noindent
\texttt{let \_ = assert (num\_digits 0 = 1)}

\pagebreak
\textit{(Extra page)}

\pagebreak
\section{Records and Variants}
Write down the type \texttt{animal} so that the following function type-checks.
\begin{lstlisting}[language=ocaml]
let get_name_and_age (a : animal) (year : int) =
  match a with
  | Cow { name = n ; age = i } -> (Some n, Some i)
  | Chicken info ->
    (Some info.name, Some (year - info.birth_year))
  | Pig age -> (None, Some age)
  | Goose -> (None, None)
\end{lstlisting}
\bigskip

\noindent\texttt{type animal =}

\pagebreak
\textit{(Extra page)}

\pagebreak
\section{List Expressions}
Circle the \texttt{list} expressions which type-check in OCaml.
\begin{enumerate}
\item \texttt{(1 :: 2 :: []) :: (3 :: [4])}
\item \texttt{((1 :: 2) :: []) :: (3 :: [])}
\item \texttt{1 :: 2 :: [3] :: 4 :: [5]}
\item \texttt{1 :: 2 :: [3] @ (4 :: [])}
\item \texttt{(1 :: 2) :: [3] @ (4 :: [])}
\end{enumerate}

\pagebreak
\textit{(Extra page)}

\pagebreak
\section{Bitonic Sequences}

A finite sequence of integers is \textbf{bitonic} if it is monotonically increasing, monotonically decreasing, or monotonically increasing and then monotonically decreasing.
That is, given $s_1, s_2, \dots, s_n$, either $s_1 < \dots < s_n$ or $s_1 > \dots  >s_n$ or there is an index $i$ such that
\begin{displaymath}
  s_1 < \dots s_{i - 1} < s_i > s_{i + 1} \dots > s_n
\end{displaymath}
Implement the function \texttt{bitonic} which, given a list of integers, returns \texttt{true} if it is bitonic, and \texttt{false} otherwise.
Your solution should be self-contained (you may write helper functions as local definitions).
\bigskip

\noindent
\texttt{let bitonic (l : int list) : bool =}
\vfill

\noindent
\texttt{let \_ = assert (bitonic [1;2;3;2;1] = true)}

\noindent
\texttt{let \_ = assert (bitonic [1;2;3] = true)}

\noindent
\texttt{let \_ = assert (bitonic [3;2;1;2] = false)}

\noindent
\texttt{let \_ = assert (bitonic [] = true)}

\noindent
\texttt{let \_ = assert (bitonic [1;1] = false)}

\noindent
\texttt{let \_ = assert (bitonic [1;2;1;2] = false)}

\pagebreak
\textit{(Extra page)}

\pagebreak
\section{Evaluation}

Consider the following pair of functions.
\begin{lstlisting}[language=ocaml]
let rec foo l =
  match l with
  | [] -> []
  | false :: bs ->
    List.map (fun x -> x - 1) (0 :: foo bs)
  | true :: bs -> bar l
and bar l =
  match l with
  | [] -> []
  | false :: bs -> foo l
  | true :: bs -> List.map ((+) 1) (0 :: bar bs)
\end{lstlisting}
\begin{enumerate}
\item
  What is the type of \texttt{foo}?
  \vspace{1in}
\item What is the value of the following expression?
  \begin{displaymath}
    \texttt{foo [true;true;true;false;false;false;false;true;true;false]}
  \end{displaymath}
\end{enumerate}

\pagebreak
\textit{(Extra page)}

\pagebreak
\section{Function Maximum}

Implement the function \texttt{func\_max} which, given \texttt{fs}, a list of functions from \texttt{int} to \texttt{int}, returns a function from \texttt{int} to \texttt{int} which is given by
\begin{displaymath}
  \mathsf{funcMax}(x) = \max(\max_{f \in \texttt{fs}}\{f(x)\}, 0)
\end{displaymath}
You should accomplish this by a single call to fold.
\begin{lstlisting}[language=ocaml]
  let op accum next = ...
  let base n = ...

  let func_max (fs : (int -> int) list) : int -> int =
    List.fold_left op base fs

  let _ = assert
    (func_max [(+) 1; fun x -> x * x] 1 = 2)
  let _ = assert
    (func_max [(+) 1; fun x -> x * x] (-2) = 4)
\end{lstlisting}
\bigskip

\noindent\texttt{let op accum next =}
\vspace{4in}

\noindent\texttt{let base n =}


\pagebreak
\textit{(Extra page)}

\pagebreak
\section{Tail Recursion}

Without using any functions from the standard library, implement a tail-recursive version of \texttt{rev\_concat}, the function which, given a list of lists, concatenates them in reverse order.
\bigskip

\noindent\texttt{rev\_concat (ls : 'a list list) : 'a list =}

\vfill

\noindent\texttt{let \_ = assert(rev\_concat [[1;2];[3;4];5] = [5;4;3;2;1])}


\pagebreak
\textit{(Extra page)}


\end{document}
